%!TEX root = ../../../thesis-trier.tex

First, in \cref{prelim:subclasses:sync:alg-unique-trap,prelim:subclasses:sync:thm:alg-unique-trap},
we investigate the running time of the problem of deciding synchronizability for automata with a trap state.
Then, \cref{thm:test_sync,prelim:thm:compute_sync_word}
make algorithmic statements for general automata. 
Lastly, in \cref{prelim:thm:comp_sync_states_if_sync},
we prove that for synchronizing automata
we can compute in linear time the set of synchronizing
states. As a consequence, stated in \cref{prelim:cor:set_sync_words_in_P},
for general automata the set of synchronizing states can be computed in quadratic time.

Later, in \cref{compl:sec:problems}, we will introduce a decision variant of the problem of computing a shortest synchronizing word.
In general, this problem is \NP-complete even for partially ordered binary automata, see~\cref{prelim:compl:sec:discussion}.

\begin{proposition}\label{prelim:subclasses:sync:alg-unique-trap}
	Let $\semiautomat{A}{Q}{\Sigma}{\delta}$.
	We can test in time $O(|Q|)$ and space $O(1)$
	whether $\mathcal A$ has a trap state, and also whether it has a unique trap state
	and compute this state if it exists.
\end{proposition}
\begin{proof}
	Iterating through all states and checking whether
	we find precisely one trap state can be done in the mentioned time and space bounds.
\end{proof} 

By \crefitem{prelim:thm:sync_trap_state}{prelim:subclasses:thm:trap:item:unique-max} 
an automaton with a trap state
is synchronizing precisely if this trap state is reachable
from every other state, which gives the following time and space bounds
to test synchronizability for automata with a trap state
that were also noted by Martyugin~\cite[Prop.~3.2]{Mar2009b}.

\begin{theorem}[{Martyugin~\cite[Prop.~3.2]{Mar2009b}}]\label{prelim:subclasses:sync:thm:alg-unique-trap}
	Let $\semiautomat{A}{Q}{\Sigma}{\delta}$ 
	be an automaton with a trap state.
	Then we can test in time $O(|Q| + |Q|\cdot |\Sigma|)$
	and space $O(|Q|)$ whether it is synchronizing.
\end{theorem} 
\begin{proof}
	By \cref{prelim:subclasses:sync:alg-unique-trap} 
	we can test whether the automaton has a unique trap, and if so compute it, in time $O(|Q|)$ and space $O(1)$.
	If this is the case, then checking whether there exists a path from every other state
	to this trap state can be done in time $O(|Q| + |\Sigma||Q|)$
	and space $O(|Q|)$ with a breadth-first search
	by following the edges in the automaton graph 
	in reverse order starting from the trap state (see, for example, the textbook~\cite{cormen01introduction} for an explanation of this algorithm
	and a statement of its resource requirements).	
\end{proof}

\begin{corollary} \label{prelim:subclasses:sync:cor:alg-unique-trap}
 For a class of automata with the property that a synchronizing state is a trap state,
 we can test whether an automaton from this class is synchronizing
 in time $O(|Q| + |Q|\cdot |\Sigma|)$
 and space $O(|Q|)$.
\end{corollary}
\begin{proof}
	First test in linear time and constant space whether an input automaton
	has a trap state, which can be done in the stated time and space bounds
	by \cref{prelim:subclasses:sync:alg-unique-trap}.
	If this is not the case, then it is not synchronizing by assumption
	Otherwise, by \cref{prelim:subclasses:sync:thm:alg-unique-trap},
	we can test in the given time and space bounds whether it is synchronizing.		
\end{proof}

By \cref{prelim:subclasses:sync:Aiffpairaut},
$\mathcal A$
is synchronizing if and only if
$\mathcal A^{\le 2} / {\sim_F}$ with $F = \{\{q\}\mid q \in Q\}$
is synchronizing.
As the pair automaton $\mathcal A^{\le 2} / {\sim_F}$ has the trap state~$F$,
$\frac{|Q|(|Q|-1)}{2} + 1$ states,
and $\left( \frac{|Q|(|Q|-1)}{2} + 1 \right) \cdot |\Sigma|$ transitions,
to test synchronizability of $\mathcal A$ we obtain the following time and space bounds from \cref{prelim:subclasses:sync:thm:alg-unique-trap}.

\begin{theorem}
	\label{thm:test_sync}
	Let $\semiautomat{A}{Q}{\Sigma}{\delta}$.
	We can test in  time $O\left(|Q|^2 + |Q|^2 \cdot |\Sigma|\right)$
	and space $O\left(|Q|^2\right)$	whether $\mathcal A$ is synchronizing.
\end{theorem}

However, the algorithm behind \cref{thm:test_sync} 
does not give us a synchronizing word. Natarajan~\cite[p.~138]{DBLP:conf/focs/Natarajan86}
has given a greedy algorithm to compute some synchronizing word of length at most $O(|Q|^3 \cdot |\Sigma|)$ in time $O(|Q|^4|\Sigma|)$.
Eppstein~\cite[Thm.~5, 6 \& 7]{Eppstein90} analyzed this algorithm further,
optimized it to run in $O(|Q|^3 + |Q|^2 \cdot |\Sigma|)$ time and $O(|Q|^2 + |Q| \cdot |\Sigma|)$ space, and proved that some synchronizing word of length at most $O(|Q|^3)$ 
can be computed by it (see also the surveys by Sandberg~\cite[Thm. 1.15]{Sandberg04}, Volkov~\cite[p.~16]{DBLP:conf/lata/Volkov08} 
and Kari \& Volkov~\cite[p.~531]{KariVolkov2021}).
Later, in \cref{compl:sec:subclasses:known-results}, we will also prove
that the deciding whether there exists a synchronizing word bounded by some given length-parameter
is \NP-complete.

\begin{theorem}[Eppstein~\cite{Eppstein90}, Natarajan~\cite{DBLP:conf/focs/Natarajan86}]
	\label{prelim:thm:compute_sync_word}
	Let $\semiautomat{A}{Q}{\Sigma}{\delta}$.
	We can test in time $O(|Q|^3 + |Q|^2 \cdot |\Sigma|)$
	and (not counting the space for the output)
	working space $O(|Q|^2 + |Q| \cdot |\Sigma|)$ whether $\mathcal A$
	is synchronizing and compute a synchronizing word of length $O(|Q|^3)$.
\end{theorem}

Now, for a synchronizing automaton, we can compute the set of synchronizing states in linear time. %time $O(|Q||\Sigma| + |Q|)$.

\begin{theorem}
	\label{prelim:thm:comp_sync_states_if_sync}
	Let $\semiautomat{A}{Q}{\Sigma}{\delta}$
	be a synchronizing automaton.
	Then we can compute in time $O(|Q| + |Q||\Sigma|)$ and space $O(|Q|)$
	the set of synchronizing states.	
\end{theorem}
\begin{proof}
	First, compute the strongly connected
	components, which can be done
	in the mentioned time and space bounds~\cite{cormen01introduction,DBLP:journals/siamcomp/Tarjan72}.
	Since~$\mathcal A$ is synchronizing,
	it follows from \cref{prelim:thm:SCC_sync_words} and \cref{prelim:prop:maximal-states-SCC}
	there exists
	a unique sink component, i.e., a strongly connected component
	with no outgoing transitions, 
	which can be easily identified and coincides with the set of synchronizing states.
	
	Alternatively, we can pick some state and start a breath-first search or a depth-first search
	in the automaton graph from this state.
	The first state $q$ we encounter with no unvisited 
	neighbors is a maximal state.
	Since~$\mathcal A$
	is synchronizing, $[q]_{\eqreach}$
	equals the set of synchronizing states
	by \cref{prelim:thm:SCC_sync_words}.
\end{proof}

For a synchronizing automaton with a trap state~$t$,
the set of synchronizing states equals $\{t\}$
and this set can be computed in linear time 
by \cref{prelim:subclasses:sync:alg-unique-trap}.

Combining Theorem~\ref{thm:test_sync} 
and~\ref{prelim:thm:comp_sync_states_if_sync}
gives the next result.

\begin{corollary}
	\label{prelim:cor:set_sync_words_in_P}
	Let $\semiautomat{A}{Q}{\Sigma}{\delta}$.
	We can compute in time 
	$
	O\left( |Q|^2 \cdot |\Sigma|\right)
	$
	and space $O\left(|Q|^2\right)$	
	the set of synchronizing states of $\mathcal A$.
\end{corollary}
\begin{proof}
	First, test whether $\mathcal A$ is synchronizing.
	If $\mathcal A$ is not synchronizing, then return the empty set.
	Otherwise, compute the set of synchronizing states and return
	this set. Combining the time and space bounds from Theorem~\ref{thm:test_sync} 
	and~\ref{prelim:thm:comp_sync_states_if_sync}
	for these computations gives the claim.
\end{proof}